\section{Ottimizzazione dell'Analisi: l'operatore split}

Questa sezione approfondisce una tecnica per accelerare l'analisi statica
riducendo la ridondanza nelle operazioni di filtraggio sui domini astratti

\subsection{Filter operator e problemi associati}

Nell'interpretazione astratta classica, la gestione del flusso di controllo
condizionale (es. \texttt{if (cond) \ldots else \ldots}) viene gestita tramite
l'operatore di \textbf{filtro}

\begin{definizione}{Filter operator}{std_filter}
L'operatore di filtro è definito come una funzione:

\[ filter: \mathbb{A} \times Pred \rightarrow \mathbb{A} \]
dove $\mathbb{A}$ è il dominio astratto e $Pred$ è un predicato (condizione).

\end{definizione}

\begin{nota}{Utilizzo dell'operatore di filtro}{filter_usage}

L'operatore di filtro (\Cref{def:std_filter}) viene utilizzato per modellare il
flusso di controllo, raffinando lo stato astratto in base alla condizione
incontrata (es. $i < 50$ nel ramo \textit{then} e $i \ge 50$ nel ramo
\textit{else}).

\begin{figure}[H]
    \centering
    \includegraphics[width=0.6\textwidth]{images/th_09/01.png}
    \caption{L'immagine illustra l'uso dell'operatore di filtro per gestire
    un ramo condizionale in un CFG. Lo stato astratto $A$ viene filtrato in
    due stati distinti in base alla condizione $c$ e alla sua negazione
    $\neg c$.}
    \label{fig:th_09_01}
\end{figure}

\end{nota}

L'approccio tradizionale presenta delle inefficienze strutturali:
\begin{enumerate}
    \item \textbf{Duplicazione del lavoro}: Per gestire un ramo condizionale,
    è necessario invocare l'operatore due volte, una per la condizione positiva
    ($filter(A, c)$) e una per la negazione ($filter(A, \neg c)$).
    \item \textbf{Overhead di memoria}: Spesso richiede la copia preliminare
    dello stato astratto $A$ per preservarlo per la seconda invocazione.
\end{enumerate}

\subsection{Split operator}

Per risolvere le inefficienze del filtro standard, viene introdotto l'operatore
\textbf{split}.

\begin{definizione}{Split operator}{split_operator}
L'operatore split combina le due operazioni di filtraggio in un'unica
invocazione atomica:
\[
split: \mathbb{A} \times Pred \rightarrow \mathbb{A} \times \mathbb{A}
\]

L'operatore prende in input uno stato astratto e un predicato, e restituisce
una coppia di stati astratti. Uno soddisfa il predicato ($P_{then}$) e uno
che soddisfa il suo complemento ($P_{else}$).
\end{definizione}

\begin{nota}{Vantaggi del Split}{split_advantages}
L'utilizzo dell'operatore split permette di:
\begin{itemize}
    \item Eliminare il lavoro replicato processando il vincolo e il suo
    opposto contemporaneamente.
    \item Ridurre gli overhead evitando copie esplicite dello stato astratto
    prima dell'operazione, delegando la gestione della memoria alla libreria
    del dominio sottostante.
\end{itemize}

\begin{figure}[H]
    \centering
    \includegraphics[width=0.6\textwidth]{images/th_09/02.png}
    \caption{L'immagine illustra l'uso dell'operatore di split per gestire
    un ramo condizionale in un CFG. Lo stato astratto $A$ viene processato
    una sola volta per produrre due stati distinti: $P_{then}$ che soddisfa la
    condizione $c$ e $P_{else}$ che soddisfa la negazione $\neg c$.}
    \label{fig:th_09_02}
\end{figure}

\end{nota}


\subsection{Implementazione e Integrazione}

L'operatore è stato implementato all'interno della libreria per poliedri PPLite
e integrato nel framework di analisi statica Crab.

\subsubsection{Modifiche al CFG}
L'integrazione richiede una modifica strutturale al Control Flow Graph (CFG)
analizzato.
Mentre un CFG standard utilizza archi annotati con \textit{assume} (o guardie)
uscenti da un blocco condizionale, la versione ottimizzata introduce uno
statement esplicito \texttt{split}.

\begin{itemize}
    \item \textbf{Standard CFG}: Archi separati con $assume(cond)$ e
    $assume(\neg cond)$.
    \item \textbf{Split CFG}: Un singolo nodo contenente l'istruzione
    $split(cond)$, che propaga internamente i due stati risultanti ai
    successori.
\end{itemize}

\begin{esempio}{Modifica del CFG con Split}{cfg_split_example}
Prendendo in considerazione il seguente frammento di codice:
\begin{verbatim}
    void foo() {
        int i = 0;
        while (i < 50) {
            i = bar(i);
        }
    }
\end{verbatim}
Da questo si possono ottenere i CFG rappresentati in \ref{fig:th_09_03}.

\begin{figure}[H]
    \centering
    \includegraphics[width=0.8\textwidth]{images/th_09/03.png}
    \caption{L'immagine illustra la differenza tra un CFG standard (a sinistra) e
    un CFG ottimizzato con l'operatore split (a destra) per il frammento di codice
    fornito. Nel CFG standard, il nodo condizionale genera due archi distinti con
    $assume(i < 50)$ e $assume(i \ge 50)$. Nel CFG con split, un singolo nodo
    esegue l'operazione $split(i < 50)$, propagando i due stati risultanti
    ai successori.}
    \label{fig:th_09_03}
\end{figure}
\end{esempio}

\subsection{Valutazione Sperimentale}

L'approccio è stato valutato su benchmark reali provenienti da progetti
differenti.
I test hanno mostrato miglioramenti significativi in termini di tempo di
analisi (\ref{fig:th_09_04}) e utilizzo della memoria (\ref{fig:th_09_05}),
senza compromettere la precisione dei risultati.


\begin{figure}[H]
    \centering
    \includegraphics[width=0.5\textwidth]{images/th_09/04.png}
    \caption{L'immagine illustra i risultati sperimentali in termini di tempo.}
    \label{fig:th_09_04}
\end{figure}


\begin{figure}[H]
    \centering
    \includegraphics[width=0.9\textwidth]{images/th_09/05.png}
    \caption{L'immagine illustra i risultati sperimentali in termini di
    utilizzo della memoria e numero di nodi del CFG.}
    \label{fig:th_09_05}
\end{figure}